%% cap03-acercamiento.tex — Nuevo acercamiento con la realidad (diseño del procedimiento)
%% Parámetros técnicos leídos de pipeline_v3/config.py. Los resultados se agregan al correr
%% el pipeline con las categorías UNESCO.

\providecommand{\pendiente}[1]{\textbf{[Pendiente: #1]}}

\chapter{Nuevo acercamiento con la realidad}
\label{cap:acercamiento}

Este capítulo describe cómo se acerca la tesis a su objeto: los documentos de política. El
procedimiento es automático de principio a fin y funciona como una sola cadena, de modo que
cada etapa usa la salida de la anterior y todo el proceso puede repetirse con los mismos
parámetros. Sigue la secuencia de la educación comparada \citep{bereday1964comparative,
adick2018bereday}: los documentos se describen y se interpretan en su contexto, se
yuxtaponen por categoría y se comparan. \textit{Este capítulo presenta el diseño; los
resultados se agregarán cuando se ejecute el procedimiento con las categorías de UNESCO.}

\section{Corpus}
\label{sec:corpus}

El corpus reúne 16 documentos oficiales de ocho países, publicados entre 2020 y 2026
(tabla~\ref{tab:corpus}). El criterio general es un documento nacional por país, emitido o
adoptado por el gobierno y vigente al 30 de junio de 2026, que trate la inteligencia
artificial y la educación. China es la excepción: aporta nueve documentos nacionales,
provinciales y municipales, porque su política se expresa en varios instrumentos. Para que
China no pese más que los demás países en las comparaciones de distancia, se excluye de
ellas el 15.º Plan Quinquenal, que por sí solo suma 181 fragmentos; con esa exclusión China
entra con 90 fragmentos y no con 271. México entra con el Plan Nacional de Inteligencia
Artificial que publicó en abril de 2026 \citep{atdt2026plan}.

\begin{table}[htbp]
\centering
\small
\caption{Documentos del corpus por país}
\label{tab:corpus}
\begin{tabular}{llllr}
\hline
País & Documento & Idioma & Año & Fragmentos \\
\hline
Alemania & Estrategia de IA del Gobierno Federal & alemán & 2020 & 230 \\
Australia & AI Action Plan & inglés & 2021 & 161 \\
Canadá & Estrategia de IA & inglés & 2026 & 210 \\
China & Nueve documentos (nacionales y locales) & chino & 2025--2026 & 271 \\
Colombia & CONPES 4144 & español & 2025 & 764 \\
Estados Unidos & AI Action Plan & inglés & 2025 & 156 \\
México & Plan Nacional de Inteligencia Artificial & español & 2026 & 116 \\
Sudáfrica & Comisión Presidencial 4IR & inglés & 2020 & 1{,}303 \\
\hline
\multicolumn{4}{l}{Total} & 3{,}211 \\
\hline
\end{tabular}
\end{table}

\section{Categorías}
\label{sec:categorias}

Las categorías se fijan antes de correr el análisis y no se ajustan después. Son dos
familias. La primera son los siete apartados de la sección 6 de la guía de UNESCO
\citep{miao2021guidance}: visión sistémica y prioridades estratégicas; principio rector;
planeación interdisciplinaria y gobernanza intersectorial; políticas y regulación para un
uso equitativo, inclusivo y ético; planes maestros para la gestión, la enseñanza, el
aprendizaje y la evaluación; pilotaje, monitoreo y evaluación; y fomento de innovaciones
locales. La segunda es la distinción de \citet{schiff2022education} entre educación para la
inteligencia artificial e inteligencia artificial para la educación. Cada categoría se
describe en un libro de códigos con su definición y los criterios de inclusión y exclusión,
tomados del texto de la guía. Un fragmento puede no pertenecer a ninguna categoría.

\section{El procedimiento}
\label{sec:procedimiento}

\subsection*{1. Fragmentación y almacenamiento}

Cada documento se conserva en su idioma original y se divide en fragmentos de 500
caracteres con un traslape de 50. Cada fragmento se guarda en una base de datos vectorial
junto con sus metadatos (país, documento, idioma, año y tipo de documento), de modo que
puede recuperarse por país. La representación vectorial se obtiene con el modelo
multilingüe \texttt{paraphrase-multilingual-MiniLM-L12-v2}, que se ejecuta de manera local
y produce vectores de 384 dimensiones comparables por similitud del coseno
\citep{reimers2019sentencebert}.

\subsection*{2. Recuperación por categoría y país}

Para cada categoría y cada país se recuperan de la base los fragmentos más cercanos a la
descripción de la categoría. La recuperación funciona como un filtro amplio: reúne
candidatos, y la decisión de si un fragmento pertenece o no a la categoría la toma el panel
en el paso siguiente. El número de candidatos por categoría y país se fija antes de la corrida completa; en la miniprueba fue de cinco.

\subsection*{3. Clasificación por el panel de modelos}

Cada fragmento candidato se entrega al panel junto con el fragmento anterior y el
posterior del mismo documento, para que el modelo lea el pasaje en su contexto y no aislado.
El panel está formado por siete modelos de lenguaje multilingües, tres de origen occidental
(\texttt{gpt-4o-mini}, \texttt{gemini-2.5-flash} y \texttt{llama-3.3-70b-instruct}) y cuatro
de origen chino (\texttt{qwen3-235b-a22b-2507}, \texttt{deepseek-chat}, \texttt{glm-4.6} y
\texttt{kimi-k2}), consultados con temperatura cero. Cada modelo decide, con el libro de
códigos, a qué categoría pertenece el fragmento y si trata de educación para la inteligencia
artificial o de inteligencia artificial para la educación. Un octavo modelo,
\texttt{claude-sonnet-4.5}, no clasifica: solo adjudica los desacuerdos del panel. Se
excluye de la clasificación porque se usó para construir el procedimiento, y hacerlo
clasificar sería circular. Usar un panel y no un solo modelo responde al riesgo de que la
elección de un modelo o de una instrucción determine el resultado
\citep{baumann2025llmhacking}.

\subsection*{4. Bases vectoriales por categoría}

Los fragmentos que el panel asigna a una categoría forman una base vectorial propia de esa
categoría, dividida por país. Al final del paso hay siete bases, una por apartado de la
guía, y en cada una los fragmentos de cada país conservan su etiqueta de educación para o
con la inteligencia artificial.

\subsection*{5. Distancias entre países dentro de cada categoría}

Dentro de cada base se calcula la distancia semántica entre los fragmentos de cada par de
países. El resultado es, para cada categoría, una matriz de distancias entre los ocho
países, que muestra cuáles quedan cerca y cuáles lejos al hablar del mismo tema. Comparar
las siete matrices responde a la segunda parte de P2: si el acomodo de los países cambia
de una categoría a otra.

\subsection*{6. Control de idioma y acuerdo del panel}

Dos pruebas ponen a prueba las distancias (P3). La primera controla el idioma: las
distancias se recalculan con todos los fragmentos traducidos a una misma lengua, y se
compara si el acomodo de los países se mantiene. Si dos países solo quedan cerca en sus
idiomas originales, su cercanía se debía al idioma y no al contenido. La segunda mide el
acuerdo entre los siete modelos del panel sobre cada fragmento, en total y separando los
modelos de origen occidental de los de origen chino. Una distancia construida sobre
fragmentos en los que el panel no se pone de acuerdo se reporta con esa advertencia.

\section{Validación}
\label{sec:validacion}

El procedimiento sigue el principio de validar cada paso \citep{grimmer2013text}. Las
categorías y las instrucciones se registran en el repositorio antes de correr el análisis,
de modo que la fecha de registro prueba que no se ajustaron a los resultados. Todas las
respuestas de los modelos se guardan tal como llegaron. Y la lectura de los resultados
vuelve a los textos: para cada distancia llamativa se examinan los fragmentos que la
producen, como pide el paso de refinamiento de \citet{nelson2020computational}.
